\section{Symmetry of the quadratic model}\label{sec:symmetry}

The remaining step is an established theorem of Vey.  We state precisely
the form we need, with the normal variable written second to match
Equation~\eqref{eq:model}.  A domain is \emph{divisible} if a group of
holomorphic automorphisms acts properly when given the discrete topology
and has a compact covering set.  This definition does not require a free
action.

\begin{theorem}[Vey]\label{thm:model-symmetry}
Let $D\subset\C^m\times\C^k$ be a connected domain biholomorphic to a
bounded domain.  Suppose that $D$ contains a point $(0,w)$ and is invariant
under
\[
 (z,w)\longmapsto(z,w+a),\qquad
 (z,w)\longmapsto(e^{i\theta}z,w),\qquad
 (z,w)\longmapsto(r^{1/2}z,rw)
\]
for every $a\in\R^k$, $\theta\in\R$, and $r>0$.
If $D$ is divisible, it is biholomorphic to a bounded symmetric domain.
The case $m=0$ is included.
\end{theorem}

In Vey's terminology these hypotheses define a divisible $S$-domain of
exponent $1/2$, after reversing the variable blocks.  His
Theorem~1 asserts homogeneity, and the immediately following consequence
asserts symmetry \cite[p.~479]{Vey1970}.  The zero-block convention is
explicit in \cite[pp.~480--481]{Vey1970}.  Neither convexity of a defining
cone nor regularity of its boundary is an input to this theorem.
Vey obtains the symmetry consequence from the earlier work of Hano,
Koszul, and Borel~\cite[p.~479]{Vey1970}; we use his combined statement rather than separately
invoke those antecedent results.

\begin{proof}[Completion of the proof of Theorem~\ref{thm:main}]
The zero-dimensional case is a point.  Otherwise,
Theorem~\ref{thm:quadratic-model} identifies $U$ biholomorphically with
the domain in Equation~\eqref{eq:model}.  We check every hypothesis of
Theorem~\ref{thm:model-symmetry}.

First we give an explicit bounded realization in the intrinsic dimension
$n=m+k$.  Boundedness of the original affine embedding places $U$ inside
$\C^N$; the realization needed here is as a domain in $\C^n$.
Extend the real forms
$\ell_1,\ldots,\ell_k$ complex linearly to $\C^k$, and write
$u_j=\ell_j(w)$.  For $(z,w)\in D$, Equation~\eqref{eq:strict-supports}
gives
\[
 \Imop u_j
 =\ell_j\bigl(\Imop w-H(z,z)\bigr)+\ell_j\bigl(H(z,z)\bigr)
 >c\norm{z}^{2}\geq0.
\]
Consequently the map
\begin{equation}\label{eq:bounded-realization}
 \mathcal B(z,w)=\left(
   \frac{u_1-i}{u_1+i},\ldots,\frac{u_k-i}{u_k+i},
   \frac{z}{u_1+i}\right)
\end{equation}
is holomorphic.  Its first $k$ coordinates have modulus less than one.
If $m>0$ and $y=\Imop u_1$, then
\[
 \norm{\frac{z}{u_1+i}}
 \leq\frac{\sqrt{y/c}}{1+y}
 \leq\frac{1}{2\sqrt c}.
\]
Here $|u_1+i|\geq1+y$, and $2\sqrt y\leq1+y$.
Thus $\mathcal B(D)$ is bounded.  Inverting the Cayley coordinates recovers
each $u_j$ holomorphically.  Independence of the $\ell_j$ then recovers
$w$, and multiplication by $u_1+i$ recovers $z$.  These formulas give a
holomorphic local inverse on an open neighborhood of every image point,
so $\mathcal B$ is a biholomorphism onto a bounded domain in $\C^{k+m}$.
When $m=0$, omit the final coordinate block and its estimate.

Choose $a_0\in C$.  Then $(0,ia_0)\in D$.  Real translations of $w$
preserve the defining residual, as do scalar unit rotations of $z$,
because $H$ is Hermitian.  Under the positive dilation in
Theorem~\ref{thm:model-symmetry}, that residual becomes
\[
 \Imop(rw)-H(r^{1/2}z,r^{1/2}z)
 =r\bigl(\Imop w-H(z,z)\bigr).
\]
Since $rC=C$ for every $r>0$, these maps preserve $D$ in both directions.

Finally, conjugate the given action on $U$ by the biholomorphism
$U\simeq D$.  Proper discontinuity and the existence of a compact covering
set are preserved by conjugacy.  This is exactly divisibility in Vey's
definition, including when point stabilizers are nontrivial.  In
particular, we do not replace the action by a free action or assume that
its quotient is a manifold.  Theorem~\ref{thm:model-symmetry} applies and
proves the assertion.
\end{proof}
